\chapter{해독: 크립에서 정지까지}\label{ch:cryptanalysis}

앞 장(\ref{ch:enigma-math}장)은 에니그마를 치환의 합성으로 적었다. 크립 위치 $k$의 평문 글자 $p$는
$c = \Stk(\Scr_k(\Stk(p)))$로 암호화된다(식~\eqref{eq:am-enigma}). 이 장은 그 식 하나에서 출발해 해독자가 키를
찾아가는 논리를 따라간다. 암호문과, 평문의 일부에 대한 짐작만 있을 때 플러그보드를 하나하나 대입하지 않고
어떻게 걸러 내는가. 끝은 봄베의 \emph{정지}와 그 뒤의 \emph{체킹}이다.

이 장에는 전선도 회로도 나오지 않는다. 나오는 것은 식과 그림이다. 단계마다 식 하나와 그림 하나를 짝지어, 같은 내용을
두 번 읽을 수 있게 했다. 작은 예는 \ref{ch:connect}장과 \ref{ch:find}장에서 다시 만날 글자 여섯 개짜리 봄베를 쓴다.
다음 두 장은 여기서 본 추론을 전선과 전압으로 다시 만든다.

\section{무차별 대입으로는 안 된다}\label{sec:ac-brute}

\Why 하루의 키는 로터 순서, 코어 위치(창 글자와 링 설정이 함께 정하는 로터의 실제 회전량), 플러그보드로 이루어진다
(\ref{sec:am-keyspace}절). 셋의 경우의 수는 크기가 전혀 다르다. Enigma~I에서 로터 순서는 60가지인데, 이 프로젝트는
반사판(B, C) 두 종까지 로터 순서에 넣어 120가지로 센다(\ref{ch:14-rust-search}장의 표와 같다). 이 120가지와 코어 위치는 $26^3 = 17{,}576$가지이므로, 둘을 합쳐 2,109,120가지다. 케이블 10개를 꽂는 플러그보드는
그 혼자서
\[
  \frac{26!}{6!\,10!\,2^{10}} = 150{,}738{,}274{,}937{,}250 \approx 1.5 \times 10^{14}
\]
가지다. 전부를 대입하려면 두 수의 곱, 약 $3.2 \times 10^{20}$가지를 시험해야 한다(그림~\ref{fig:ac-keyspace}).

\begin{figure}[ht]
\centering
\begin{tikzpicture}[x=4.5mm,y=1cm]
  \foreach \e in {0,2,...,20} {
    \draw[cgray!35] (\e,0.35) -- (\e,-2.45);
    \node[note,font=\scriptsize] at (\e,-2.7) {$10^{\e}$};
  }
  \fill[cblue] (0,-0.18) rectangle (6.32,0.18);
  \node[anchor=east,font=\footnotesize] at (-0.4,0) {로터 순서 $\times$ 코어 위치};
  \node[anchor=west,font=\footnotesize] at (6.5,0) {2,109,120};
  \fill[corange] (0,-1.18) rectangle (14.18,-0.82);
  \node[anchor=east,font=\footnotesize] at (-0.4,-1) {플러그보드(10쌍)};
  \node[anchor=west,font=\footnotesize] at (14.4,-1) {$1.5\times10^{14}$};
  \fill[cred] (0,-2.18) rectangle (20.5,-1.82);
  \node[anchor=east,font=\footnotesize] at (-0.4,-2) {전부 대입};
  \node[anchor=west,font=\footnotesize] at (20.7,-2) {$3.2\times10^{20}$};
\end{tikzpicture}
\caption{Enigma~I의 키를 나눈 크기(가로축은 로그 눈금). 봄베가 훑는 것은 파란 막대뿐이고, 주황 막대는 논리로 걸러 낸다.}
\label{fig:ac-keyspace}
\end{figure}

시간으로 바꿔 보면 차이가 실감 난다. \ref{sec:am-keyspace}절은 초당 10억 번이라는 넉넉한 가정으로 어림했다. 여기서는
이 프로젝트의 실제 속도로 다시 본다. Rust 모델은 PC 한 대에서 초당 약 1,000만 위치를 시험한다(보고서 7장).
파란 막대 전체는 0.2초쯤이다. 같은 빠르기로 빨간 막대를 시험한다면 약 $3 \times 10^{13}$초, 곧 100만 년쯤
걸린다. 1940년대의 기계는 이보다 훨씬 느렸다.

그래서 해독의 첫 착상은 \emph{키를 둘로 가르는 것}이다. 로터 순서와 코어 위치는 수가 적으니 기계로 하나씩 훑는다.
플러그보드는 대입하지 않고, 위치마다 ``이 위치에서 말이 되는 플러그보드가 있을 수 있는가''를 논리로 따진다.
이 장의 나머지는 그 논리가 어디서 오는지를 본다.

\Ours 일의 양을 실제로 세는 코드는 \ref{sec:s-count}절에 있다. 그 표의 ``전체 위치''가 파란 막대다.

\section{크립과 정렬}\label{sec:ac-crib}

논리에는 재료가 필요하다. 그 재료가 \term{크립}{crib}, 곧 평문에 들어 있으리라고 짐작하는 문구다. 독일군 전문은
형식이 정해져 있어 이런 짐작이 가능했다. 예로 \texttt{WETTERVORHERSAGE}(일기예보, 16글자)를 쓴다. 이 예의
암호문(121글자)과 결과는 모두 보고서 3장과 7장의 것이다.

크립이 암호문의 어디에 놓이는지는 모른다. 하지만 놓을 수 없는 자리는 안다. 에니그마는 어떤 글자도 자기 자신으로
바꾸지 않는다(\ref{sec:am-properties}절). 크립을 시작 자리 $o$에 놓으면 크립의 $k$번째 글자 $p_k$ 아래에 암호문 글자
$c_{o+k}$가 온다. 둘이 같은 자리가 하나라도 있으면 그 $o$는 불가능하다(그림~\ref{fig:ac-align}).

\begin{figure}[ht]
\centering
\begin{tikzpicture}[x=4.6mm,y=5.5mm,font=\ttfamily\small]
  \node[anchor=east,font=\small] at (-0.8,0) {암호문};
  \node[anchor=east,font=\small] at (-0.8,-1) {크립};
  \foreach \c [count=\i from 0] in {F,L,O,G,G,S,L,V,I,X,B,T,V,W,V,E} { \node at (\i,0) {\c}; }
  \foreach \c [count=\i from 0] in {W,E,T,T,E,R,V,O,R,H,E,R,S,A,G,E} { \node at (\i,-1) {\c}; }
  \draw[cred,line width=0.9pt,rounded corners=1pt] (14.5,0.5) rectangle (15.5,-1.5);
  \node[anchor=west,font=\small,text=cred] at (16.2,-0.5) {자리 2: 불가능 (\texttt{E} 아래 \texttt{E})};
  \begin{scope}[yshift=-17mm]
  \node[anchor=east,font=\small] at (-0.8,0) {암호문};
  \node[anchor=east,font=\small] at (-0.8,-1) {크립};
  \foreach \c [count=\i from 0] in {I,L,Z,L,I,I,U,F,M,W,A,J,N,W,T,U} { \node at (\i,0) {\c}; }
  \foreach \c [count=\i from 0] in {W,E,T,T,E,R,V,O,R,H,E,R,S,A,G,E} { \node at (\i,-1) {\c}; }
  \foreach \i in {1,...,16} { \node[font=\tiny\rmfamily,text=cgray] at (\i-1,-1.8) {\i}; }
  \node[anchor=west,font=\small,text=cgreen] at (16.2,-0.5) {자리 28: 가능 (실제 자리)};
  \end{scope}
\end{tikzpicture}
\caption{크립 정렬(보고서 3장의 예). 같은 글자가 위아래로 겹치는 자리는 배제된다. 아래 줄의 작은 숫자는 크립 안의
위치 $k$다. 자리 28에는 겹치는 글자가 없다.}
\label{fig:ac-align}
\end{figure}

\Why 이 검사가 얼마나 걸러 내는지 어림할 수 있다. 틀린 자리에서는 위아래 글자가 서로 무관하다고 보면, 한 자리에서
우연히 겹치지 않을 확률은 $25/26$이다. 16글자가 모두 겹치지 않아야 살아남으므로
\begin{equation}
  \Pr(\text{틀린 자리가 살아남음}) \approx \left(\frac{25}{26}\right)^{16} \approx 0.53 . \label{eq:ac-survive}
\end{equation}
예제 전문에서 크립을 놓을 수 있는 자리는 $121 - 16 + 1 = 106$곳이고, 이 검사로 46곳이 지워져 60곳이 남는다(보고서 3장).
어림으로는 틀린 자리 105곳 가운데 약 56곳이 살아남고 실제 자리 하나를 더해 약 57곳이니, 실제 60곳과 비슷하다.

이 검사는 싸지만 약하다. 자리의 절반쯤만 지운다. 크립이 길수록 더 지우지만, 정렬만으로 답이 나오지는 않는다.
남은 자리마다 아래의 논리를 다시 적용한다.

\Ours \code{possible\_offsets}(\ref{ch:rust-menu}장)가 이 검사를 한 줄로 한다.

\section{플러그보드를 한쪽으로 넘기기}\label{sec:ac-move}

크립을 한 자리에 놓으면 ``크립 위치 $k$에서 $p$가 $c$가 되었다''는 사실이 글자 수만큼 생긴다. 사실 하나는
식~\eqref{eq:am-enigma}다.
\[
  c = \Stk\bigl(\Scr_k(\Stk(p))\bigr).
\]
이 식에는 모르는 것이 둘 섞여 있다. $\Scr_k$는 로터 순서와 코어 위치를 정하면 계산할 수 있다. 봄베가 하나씩
훑을 것이므로 ``안다''고 놓는다. $\Stk$는 끝까지 모른다. 그런데 $\Stk$가 양쪽 끝에 하나씩 있다.

\Why 플러그보드는 쌍을 맞바꾸므로 두 번 지나면 제자리다. $\Stk(\Stk(x)) = x$다. 식의 양변에 $\Stk$를 씌우면
오른쪽 바깥의 $\Stk$가 지워진다.
\begin{equation}
  \Stk(c) = \Scr_k\bigl(\Stk(p)\bigr). \label{eq:ac-moved}
\end{equation}
모르는 것이 한쪽으로 모였다. 남은 식은 ``$p$의 짝을 스크램블러에 넣으면 $c$의 짝이 나온다''이다
(그림~\ref{fig:ac-path}, \ref{fig:ac-square}).

\begin{figure}[ht]
\centering
\begin{tikzpicture}[x=1cm,y=1cm,
    unk/.style={draw=corange,fill=corangebg,minimum size=7mm,rounded corners=1pt,font=\small,inner sep=2pt}]
  \node[ltr] (p) at (0,0) {$p$};
  \node[blko,minimum width=10mm] (s1) at (1.7,0) {$\Stk$};
  \node[unk] (u1) at (3.4,0) {?};
  \node[blkb,minimum width=12mm] (sc) at (5.2,0) {$\Scr_k$};
  \node[unk] (u2) at (7.0,0) {?};
  \node[blko,minimum width=10mm] (s2) at (8.7,0) {$\Stk$};
  \node[ltr] (c) at (10.4,0) {$c$};
  \draw[->] (p) -- (s1); \draw[->] (s1) -- (u1); \draw[->] (u1) -- (sc);
  \draw[->] (sc) -- (u2); \draw[->] (u2) -- (s2); \draw[->] (s2) -- (c);
  \node[note,above=1mm of u1] {$\Stk(p)$};
  \node[note,above=1mm of u2] {$\Stk(c)$};
  \node[note,anchor=east] at (-0.5,0) {(a)};
  \node[unk] (v1) at (3.4,-1.8) {$x$};
  \node[blkb,minimum width=12mm] (sc2) at (5.2,-1.8) {$\Scr_k$};
  \node[unk,anchor=west] (v2) at (6.65,-1.8) {$\Scr_k(x)$};
  \draw[->] (v1) -- (sc2); \draw[->] (sc2) -- (v2);
  \draw[cgray,dashed] (u1) -- (v1); \draw[cgray,dashed] (u2) -- (u2 |- v2.north);
  \node[note,below=1mm of v1] {가설 $\Stk(p)=x$};
  \node[note,below=1mm of v2] {귀결 $\Stk(c)=\Scr_k(x)$};
  \node[note,anchor=east] at (-0.5,-1.8) {(b)};
\end{tikzpicture}
\caption{(a) 에니그마 안의 전류의 길. 주황 칸 둘은 모르는 글자다. 오른쪽 칸이 $\Stk(c)$인 것은 거기에 $\Stk$를
씌우면 $c$가 되기 때문이다. (b) 양 끝의 $\Stk$를 떼어 낸 꼴, 식~\eqref{eq:ac-moved}. 주황 칸 하나를 짐작하면 다른
하나는 계산된다.}
\label{fig:ac-path}
\end{figure}

\begin{figure}[ht]
\centering
\begin{tikzpicture}[x=1cm,y=1cm]
  \node (p) at (0,0) {$p$};
  \node (c) at (4.6,0) {$c$};
  \node (Sp) at (0,-2.2) {$\Stk(p)$};
  \node (Sc) at (4.6,-2.2) {$\Stk(c)$};
  \draw[->,cgray] (p) -- node[above,note] {크립 위치 $k$의 에니그마} (c);
  \draw[<->,corange,line width=0.9pt] (p) -- node[left,text=corange] {$\Stk$} (Sp);
  \draw[<->,corange,line width=0.9pt] (c) -- node[right,text=corange] {$\Stk$} (Sc);
  \draw[->,cblue,line width=1.1pt] (Sp) -- node[below,text=cblue] {$\Scr_k$} (Sc);
  \node[note,anchor=west,align=left] at (6.0,-1.1) {위: 해독자가 보는 것 ($p$와 $c$)\\
     아래: 해독자가 계산할 수 있는 것 ($\Scr_k$)\\
     세로: 모르는 것 ($\Stk$, 양방향)};
\end{tikzpicture}
\caption{같은 내용을 네모로 그린 것. 전류는 왼쪽 아래로, 아래쪽을 건너, 오른쪽 위로 간다. 세로 화살표가 양쪽을
향하는 것은 $\Stk$를 거꾸로 지나도 $\Stk$이기 때문이다. 그래서 오른쪽 아래 모서리는 ``$\Stk(p)$를 $\Scr_k$로 보낸
것''이자 ``$c$를 $\Stk$로 보낸 것''이다. 이것이 식~\eqref{eq:ac-moved}다.}
\label{fig:ac-square}
\end{figure}

식~\eqref{eq:ac-moved}가 주는 것은 \emph{추론}이다. ``$\Stk(p) = x$라면 $\Stk(c) = \Scr_k(x)$다.'' 앞으로
``글자 $X$의 짝은 $y$다'', 곧 $\Stk(X) = y$라는 진술을 \term{가설}{hypothesis}이라 부른다. 글자가 26개이고 짝의
후보가 26개이므로 가설은 $26 \times 26 = 676$개다. 식~\eqref{eq:ac-moved}는 가설 하나에서 다른 가설 하나를 끌어내는
규칙이다.

이 규칙은 거꾸로도 쓸 수 있다. $\Scr_k$도 되돌릴 수 있는 치환이어서($\Scr_k(\Scr_k(y)) = y$,
\ref{sec:am-properties}절) 식~\eqref{eq:ac-moved}의 양변에 $\Scr_k$를 씌우면 $\Stk(p) = \Scr_k(\Stk(c))$가 된다.
``$\Stk(c) = y$라면 $\Stk(p) = \Scr_k(y)$다.'' 평문 쪽에서 출발하든 암호문 쪽에서 출발하든 같은 규칙이다. 이 대칭이
뒤에서 여러 번 쓰인다.

\General 두 사람의 나이 차가 3살인 것을 알면, 한 사람의 나이를 짐작하는 순간 다른 사람의 나이가 정해진다. 어느
쪽 나이를 먼저 짐작해도 된다. 식~\eqref{eq:ac-moved}는 ``나이 차'' 대신 치환 $\Scr_k$가 두 글자의 짝을 묶는다.

\Ours \ref{ch:find}장의 작은 봄베에서 크립 위치 1의 스크램블러는 $\Scr_1 = \mathtt{a}\leftrightarrow\mathtt{b}$,
$\mathtt{c}\leftrightarrow\mathtt{f}$, $\mathtt{d}\leftrightarrow\mathtt{e}$이고, 크립 위치 1에서 \texttt{A}가 \texttt{B}가
되었다. 그러면 ``$\Stk(\mathtt{A}) = \mathtt{d}$라면 $\Stk(\mathtt{B}) = \Scr_1(\mathtt{d}) = \mathtt{e}$''다. 봄베는 이
추론 하나를 전선 한 가닥으로 만든다. 그림~\ref{fig:c-scr}(\ref{sec:c-scrambler}절)의 파란 선이 바로 이 추론이다.
Rust 모델의 체킹 단계도 이 식을 그대로 쓴다(\ref{ch:rust-check}장의 \code{extend\_with\_crib}).

\section{메뉴}\label{sec:ac-menu}

크립 16글자는 식~\eqref{eq:ac-moved}를 16개 준다. 이것을 한눈에 보려고 그래프를 그린다. 글자를 점으로 놓고, 크립
위치 $k$에서 $p$가 $c$가 되었으면 두 점을 선으로 잇고 선에 $k$를 적는다. 이 그래프가 \term{메뉴}{menu}다.
선에는 방향이 없다. 앞 절의 끝에서 본 대로 식~\eqref{eq:ac-moved}는 양쪽으로 읽히기 때문이다
(그림~\ref{fig:ac-edge}).

\begin{figure}[ht]
\centering
\begin{tikzpicture}[x=1cm,y=1cm]
  \node[ltr] (E) at (0,0) {E};
  \node[ltr] (I) at (3,0) {I};
  \draw (E) -- node[above,font=\scriptsize,text=cgray] {5} (I);
  \node[note,anchor=north] at (1.5,-0.45) {크립 위치 5: 평문 \texttt{E}, 암호문 \texttt{I}};
  \node[anchor=west,align=left,font=\small] (eq) at (4.2,0) {$\Stk(\mathtt{I}) = \Scr_5\bigl(\Stk(\mathtt{E})\bigr)$\\[3pt]
    $\Stk(\mathtt{E}) = \Scr_5\bigl(\Stk(\mathtt{I})\bigr)$};
  \draw[decorate,decoration={brace,amplitude=4pt}] ([xshift=1mm]eq.north east) -- ([xshift=1mm]eq.south east)
    node[midway,right=5pt,note,align=left] {같은 식\\(어느 끝에서든)};
\end{tikzpicture}
\caption{메뉴의 선 하나. 그림~\ref{fig:ac-align}의 자리 28에서 크립의 다섯째 글자 \texttt{E} 아래에 \texttt{I}가 있다.
선은 식~\eqref{eq:ac-moved} 하나이고, 방향이 없다.}
\label{fig:ac-edge}
\end{figure}

그림~\ref{fig:g-menu}는 자리 28에 놓은 크립 전체로 만든 메뉴다.

\begin{figure}[ht]
\centering
\begin{tikzpicture}[x=1cm,y=1cm,
    ed/.style={line width=0.6pt},
    lp/.style={line width=1.4pt,draw=cblue},
    el/.style={font=\scriptsize,fill=white,inner sep=1pt,text=cgray}]
  \node[ltr,fill=corangebg,draw=corange,line width=1pt] (E) at (0,0) {E};
  \node[ltr] (I) at (2,1.2) {I};
  \node[ltr] (W) at (4,0) {W};
  \node[ltr] (A) at (2,-1.2) {A};
  \node[ltr] (L) at (-2,1.2) {L};
  \node[ltr] (T) at (-4,1.2) {T};
  \node[ltr] (Z) at (-6,1.9) {Z};
  \node[ltr] (G) at (-6,0.5) {G};
  \node[ltr] (U) at (-2,-1.2) {U};
  \node[ltr] (V) at (-4,-1.2) {V};
  \node[ltr] (R) at (4,2.3) {R};
  \node[ltr] (M) at (6,2.9) {M};
  \node[ltr] (J) at (6,1.7) {J};
  \node[ltr] (H) at (6,0) {H};
  \draw[lp] (E) -- node[el] {5} (I);
  \draw[lp] (I) -- node[el] {1} (W);
  \draw[lp] (W) -- node[el] {14} (A);
  \draw[lp] (A) -- node[el] {11} (E);
  \draw[ed] (E) -- node[el] {2} (L);
  \draw[ed] (L) -- node[el] {4} (T);
  \draw[ed] (T) -- node[el] {3} (Z);
  \draw[ed] (T) -- node[el] {15} (G);
  \draw[ed] (E) -- node[el] {16} (U);
  \draw[ed] (U) -- node[el] {7} (V);
  \draw[ed] (I) -- node[el] {6} (R);
  \draw[ed] (R) -- node[el] {9} (M);
  \draw[ed] (R) -- node[el] {12} (J);
  \draw[ed] (W) -- node[el] {10} (H);
  \node[note,text=corange,anchor=east,font=\scriptsize] at (-0.42,0) {테스트 글자};
  \node[note,text=cblue] at (2,0) {루프};
\end{tikzpicture}
\caption{그림~\ref{fig:ac-align}의 자리 28에서 만든 메뉴(보고서 3장). 글자 14개, 선 14개, 루프 1개(파랑).
선의 숫자는 크립 위치 $k$다. 크립 16글자 가운데 이 덩어리에 이어지지 않는 두 관계(\texttt{O}--\texttt{F},
\texttt{S}--\texttt{N})는 뺐다.}
\label{fig:g-menu}
\end{figure}

메뉴를 보면 쓸모 있는 부분과 없는 부분이 갈린다.

\begin{itemize}
\item \textbf{덩어리.} 추론은 선을 따라서만 간다. 그러므로 서로 이어진 글자들의 \term{덩어리}{connected component}
하나 안에서만 가설이 서로 영향을 준다. \texttt{O}--\texttt{F} 선은 \texttt{E}의 짝에 대해 아무것도 말해 주지 않는다.
봄베는 가장 큰 덩어리를 쓴다.
\item \textbf{루프.} \texttt{E}--\texttt{I}--\texttt{W}--\texttt{A}--\texttt{E}처럼 출발한 글자로 돌아오는 길이
\term{루프}{loop}다. 다음 절에서 보듯 걸러 내는 힘은 루프에서 나온다. 덩어리 하나의 루프 수는
(선 수) $-$ (글자 수) $+$ 1이다. 이 메뉴는 $14 - 14 + 1 = 1$이다.
\end{itemize}

같은 전문에서 크립을 23글자(\texttt{WETTERVORHERSAGEFUERDIE})로 늘리고 선 14개를 골라 만든 메뉴는 글자 11개, 루프
4개다(\ref{ch:rust-menu}장). 선의 수(곧 필요한 스크램블러의 수)는 같은데 루프가 넷이다. 이 차이가 어떤 결과를 낳는지는
\ref{sec:ac-stop}절에서 본다.

\Ours 메뉴를 만들고 루프를 세고 선을 고르는 코드가 \ref{ch:rust-menu}장이다. 루프 수의 공식이 왜 그런지도 그 장의
그림~\ref{fig:g-loops}에 있다.

\section{루프의 조건}\label{sec:ac-loop}

\Why 메뉴의 길 하나를 따라가 보자. 글자 $X_0$에서 크립 위치 $k_1$의 선으로 $X_1$에 가고, $k_2$의 선으로 $X_2$에
가고, \ldots, $k_n$의 선으로 $X_n$에 간다. 선마다 식~\eqref{eq:ac-moved}를 쓰면(방향은 상관없다)
\[
  \Stk(X_1) = \Scr_{k_1}\bigl(\Stk(X_0)\bigr),\quad
  \Stk(X_2) = \Scr_{k_2}\bigl(\Stk(X_1)\bigr),\quad \ldots
\]
이고, 이것을 차례로 대입하면 $\Stk(X_n) = \Scr_{k_n} \cdots \Scr_{k_2}\Scr_{k_1}\bigl(\Stk(X_0)\bigr)$다. 치환의 합성은
오른쪽부터 적용한다(\ref{ch:enigma-math}장). 길이 출발점으로 돌아오면, 곧 $X_n = X_0$이면 이렇게 된다.

\begin{keybox}[title={루프의 조건}]
루프 $X_0 \to X_1 \to \cdots \to X_n = X_0$의 선이 크립 위치 $k_1, \ldots, k_n$이면, 맞는 로터 순서와 코어 위치에서
\begin{equation}
  \Stk(X_0) = P\bigl(\Stk(X_0)\bigr), \qquad P = \Scr_{k_n} \cdots \Scr_{k_2}\, \Scr_{k_1}. \label{eq:ac-loop}
\end{equation}
즉 $\Stk(X_0)$는 스크램블러들의 곱 $P$의 \term{고정점}{fixed point}이어야 한다.
\end{keybox}

$P$는 위치만 정하면 계산되는 치환이다. 위치마다 $P$를 계산해 고정점을 보면, $X_0$의 짝이 될 수 있는 글자가 고정점들로
좁혀진다. 고정점이 하나도 없으면 이 위치에서는 메뉴를 만족하는 플러그보드가 없다.

얼마나 좁혀지는가. 틀린 위치에서 $P$를 아무렇게나 섞인 치환으로 보면, 글자 하나가 $P$로 자기 자신에 갈 확률은
$1/26$이다. 그래서 가설 26개 가운데 루프를 통과하는 것은 평균 하나다(무작위 순열의 고정점 수의 평균은 정확히 1이다).
루프 하나가 후보를 약 26분의 1로 줄인다. 루프가 여럿이면 짝은 모든 루프의 $P$에 공통인 고정점이어야 하므로 루프마다
다시 26분의 1이 곱해진다.

\Ours 작은 봄베(\ref{ch:find}장)의 메뉴는 루프 하나, \texttt{A}--\texttt{B}--\texttt{C}--\texttt{A}다. 크립 위치 1, 2, 3의
선에 스크램블러 $\Scr_1$, $\Scr_2$, $\Scr_3$이 걸린다. \ref{ch:connect}장과 \ref{ch:find}장의 그림도 같은 이름을 쓴다. 글자가 여섯뿐이므로 $P = \Scr_3\Scr_2\Scr_1$을 손으로 계산할 수 있다. 드럼 위치 1에서
\[
  \Scr_1 = (\mathtt{ab})(\mathtt{cf})(\mathtt{de}),\quad \Scr_2 = (\mathtt{ab})(\mathtt{cd})(\mathtt{ef}),\quad
  \Scr_3 = (\mathtt{ac})(\mathtt{be})(\mathtt{df})
\]
이다. $(\mathtt{ab})$는 \texttt{a}와 \texttt{b}를 맞바꾼다는 뜻이다. \texttt{a}를 넣으면 $\Scr_1$이 \texttt{b}로,
$\Scr_2$가 \texttt{a}로, $\Scr_3$이 \texttt{c}로 보내므로 $P(\mathtt{a}) = \mathtt{c}$다. 여섯 글자를 모두 하면
\[
  P = (\mathtt{a}\,\mathtt{c}\,\mathtt{b}\,\mathtt{e})\,(\mathtt{d})\,(\mathtt{f})
\]
이다. \texttt{a}$\to$\texttt{c}$\to$\texttt{b}$\to$\texttt{e}$\to$\texttt{a}로 도는 순환 하나와 고정점 \texttt{d},
\texttt{f}다. 그러므로 $\Stk(\mathtt{A})$는 \texttt{d}나 \texttt{f}만 될 수 있다. 드럼 위치 2에서는 $\Scr_3$만
$(\mathtt{ac})(\mathtt{bd})(\mathtt{ef})$로 바뀌고,
\[
  P = (\mathtt{a}\,\mathtt{c}\,\mathtt{f}\,\mathtt{b}\,\mathtt{d}\,\mathtt{e})
\]
로 여섯 글자가 한 순환을 이룬다. 고정점이 없으므로 위치 2는 버린다(그림~\ref{fig:ac-loop}).

\begin{figure}[ht]
\centering
\begin{tikzpicture}[x=1cm,y=1cm,
    el/.style={font=\scriptsize,fill=white,inner sep=1pt,text=cgray},
    fx/.style={ltr,fill=cgreenbg,draw=cgreen,line width=0.9pt}]
  % 메뉴
  \node[ltr] (A) at (0,0) {A};
  \node[ltr] (B) at (2,0) {B};
  \node[ltr] (C) at (1,1.6) {C};
  \draw[line width=1.2pt,cblue] (A) -- node[el] {1} (B) -- node[el] {2} (C) -- node[el] {3} (A);
  \node[note] at (1,-0.8) {메뉴};
  % 드럼 위치 1
  \begin{scope}[xshift=4.2cm]
    \node[ltr] (a) at (0,1.4) {a};
    \node[ltr] (c) at (1.3,1.4) {c};
    \node[ltr] (b) at (1.3,0.1) {b};
    \node[ltr] (e) at (0,0.1) {e};
    \draw[->] (a) -- (c); \draw[->] (c) -- (b); \draw[->] (b) -- (e); \draw[->] (e) -- (a);
    \node[fx] (d) at (2.5,0.1) {d};
    \node[fx] (f) at (3.4,0.1) {f};
    \path (d) edge[->,loop,out=130,in=50,min distance=6mm,cgreen,line width=0.8pt] (d);
    \path (f) edge[->,loop,out=130,in=50,min distance=6mm,cgreen,line width=0.8pt] (f);
    \node[note] at (1.7,-0.8) {드럼 위치 1: 고정점 \texttt{d}, \texttt{f}};
  \end{scope}
  % 드럼 위치 2
  \begin{scope}[xshift=10.4cm,yshift=0.75cm]
    \foreach \l/\ang in {a/90,c/30,f/-30,b/-90,d/-150,e/150} { \node[ltr] (q\l) at (\ang:1.05) {\l}; }
    \draw[->] (qa) -- (qc); \draw[->] (qc) -- (qf); \draw[->] (qf) -- (qb);
    \draw[->] (qb) -- (qd); \draw[->] (qd) -- (qe); \draw[->] (qe) -- (qa);
    \node[note] at (0,-1.55) {드럼 위치 2: 고정점 없음};
  \end{scope}
\end{tikzpicture}
\caption{왼쪽은 작은 봄베의 메뉴, 오른쪽 둘은 같은 루프의 곱을 화살표로 그린 것이다(${P = \Scr_3\Scr_2\Scr_1}$).
화살표 하나가 ``루프를 한 바퀴 돌면 이 글자가 저 글자가 된다''이다. 초록 글자는 한 바퀴 돌아 제자리로 오는 고정점이고,
$\Stk(\mathtt{A})$가 될 수 있는 것은 이것뿐이다.}
\label{fig:ac-loop}
\end{figure}

\paragraph{가설 하나로 모두를 시험한다.} 고정점을 찾으려고 26글자를 하나씩 넣어 볼 필요는 없다. 아무 가설이나
하나를 넣고 귀결을 끝까지 따라가면 된다. 드럼 위치 1에서 $\Stk(\mathtt{A}) = \mathtt{a}$를 넣고 루프를 계속 돌면
그림~\ref{fig:ac-orbit}이 된다. 한 바퀴마다 $\Stk(\mathtt{A})$의 새 값이 나오고, 네 바퀴 만에 처음의 \texttt{a}로
돌아온다. 나온 값은 \texttt{a}, \texttt{c}, \texttt{b}, \texttt{e}, 곧 $P$에서 \texttt{a}가 속한 순환이다.

\begin{figure}[ht]
\centering
\begin{tikzpicture}[x=2.0cm,y=-9mm,
    lb/.style={font=\scriptsize,text=cgray,above}]
  \node[font=\small] at (0,-0.95) {$\Stk(\mathtt{A})$};
  \node[font=\small] at (1,-0.95) {$\Stk(\mathtt{B})$};
  \node[font=\small] at (2,-0.95) {$\Stk(\mathtt{C})$};
  \node[font=\small] at (3,-0.95) {$\Stk(\mathtt{A})$};
  \foreach \r/\x/\y/\z/\w in {0/a/b/a/c,1/c/f/e/b,2/b/a/b/e,3/e/d/c/a} {
    \node[ltr,draw=cred,line width=0.9pt] (n\r0) at (0,\r) {\x};
    \node[ltr] (n\r1) at (1,\r) {\y};
    \node[ltr] (n\r2) at (2,\r) {\z};
    \node[ltr] (n\r3) at (3,\r) {\w};
    \draw[->] (n\r0) -- node[lb] {$\Scr_1$} (n\r1);
    \draw[->] (n\r1) -- node[lb] {$\Scr_2$} (n\r2);
    \draw[->] (n\r2) -- node[lb] {$\Scr_3$} (n\r3);
    \ifnum\r<3 \node[note,anchor=west] at (3.25,\r) {$\Rightarrow$ 다음 줄의 처음}; \fi
  }
  \node[note,anchor=west,text=cgreen] at (3.25,3) {$\Rightarrow$ 처음의 \texttt{a}로 돌아옴};
  \node[note,text=cred,anchor=north,align=center] at (0,3.55) {\texttt{A}의 짝이 넷:\\\texttt{a}, \texttt{c}, \texttt{b}, \texttt{e}};
\end{tikzpicture}
\caption{드럼 위치 1에서 가설 $\Stk(\mathtt{A}) = \mathtt{a}$를 넣고 루프를 네 바퀴 따라간 것. 줄 하나가 한 바퀴이고,
줄 끝의 값이 다음 줄의 처음이 된다. 빨간 칸이 \texttt{A}의 짝으로 나온 값이다.}
\label{fig:ac-orbit}
\end{figure}

이 넷은 서로에게서 끌어낸 것이다. 그리고 앞 절에서 본 대로 추론은 거꾸로도 성립하므로, 넷 가운데 어느 것에서
출발해도 나머지 셋이 나온다. 그러니 넷은 함께 참이거나 함께 거짓이다. 함께 참일 수는 없다. 한 글자의 짝은 하나뿐이다.
그러므로 넷은 모두 거짓이다. 거짓일지도 모르는 가설 하나를 넣었을 뿐인데, 가설 넷이 한꺼번에 지워졌다. 남은 것은
순환 밖에 있던 \texttt{d}와 \texttt{f}, 곧 고정점이다.

일반적으로 적으면 이렇다. 가설 $\Stk(X_0) = x$에서 출발한 귀결은 $x, P(x), P^2(x), \ldots$이고, 이것은 $P$의 순환
하나다. 순환의 길이가 2 이상이면 그 순환의 가설은 모두 거짓이고, 길이가 1이면(고정점) 살아남는다. 루프가 여럿이거나
루프에 다른 가지가 붙은 메뉴에서도 같다. 가설 하나에서 끌어낼 수 있는 것을 모두 모은 \emph{무리}는 함께 서고 함께
넘어진다. 테스트 글자에 대한 가설 26개가 모두 어느 무리에선가 지워지면, 이 위치에는 메뉴를 만족하는 플러그보드가
없다. 루프가 테스트 글자를 지나지 않아도 된다. 같은 덩어리에 있으면 테스트 글자에서 그 루프까지 가는 길을 따라
가설이 옮겨 가고, 같은 길로 돌아오기 때문이다.

튜링이 설계한 봄베~\cite{jackson-bh06,ellsbury1998bombe}가 위치마다 하는 일이 이것이다. 가설 하나를 넣고 귀결을
퍼뜨린 다음, 테스트 글자의 26개 가설이 모두 지워졌는지를 본다. 넣은 가설이 맞든 틀리든 상관없다.

\Ours 그림~\ref{fig:ac-orbit}에서 나온 값을 글자별로 모으면 \texttt{A}는 \texttt{a}, \texttt{b}, \texttt{c}, \texttt{e},
\texttt{B}는 \texttt{a}, \texttt{b}, \texttt{d}, \texttt{f}, \texttt{C}는 \texttt{a}, \texttt{b}, \texttt{c}, \texttt{e}다. \ref{ch:find}장의 그림~\ref{fig:f-groups}에서 전압을 건 선과
함께 켜지는 파란 덩어리 12가닥이 정확히 이것이다. 루프를 ``돌리는'' 것을 전류가 한꺼번에 하는 것이다.

\section{대각선 보드: 짝은 서로를 가리킨다}\label{sec:ac-diagonal}

지금까지 쓴 사실은 식~\eqref{eq:ac-moved}와 ``짝은 하나뿐''이었다. 플러그보드에 대해 아직 쓰지 않은 사실이 하나 있다.
케이블 하나가 두 글자를 맞바꾸므로
\begin{equation}
  \Stk(X) = y \iff \Stk(Y) = x \label{eq:ac-sym}
\end{equation}
이다. ``$X$의 짝이 $y$''와 ``$Y$의 짝이 $x$''는 같은 말이다. 글자와 짝을 맞바꾼 \emph{거울쌍}이다
(그림~\ref{fig:ac-mirror}의 위). 이것도 추론 규칙이다. 그런데 식~\eqref{eq:ac-moved}와 성격이 다르다. 메뉴의 선을
따라가지 않고, 한 글자에 대한 가설을 \emph{다른 글자}에 대한 가설로 옮긴다. 귀결에 메뉴의 다른 글자가 짝으로 나오면
그 글자의 가설로 건너가 다시 퍼진다. 루프가 하던 일, 곧 추론을 출발한 곳으로 되돌려 보내는 일을 루프가 없는 곳에서도
한다.

\begin{figure}[ht]
\centering
\begin{tikzpicture}[x=1cm,y=1cm,
    hy/.style={draw,rounded corners=1.5pt,fill=white,minimum height=7mm,minimum width=19mm,font=\small,inner sep=2pt},
    lb/.style={font=\scriptsize,text=cgray,above}]
  % (a) 거울쌍
  \node[hy,draw=corange,fill=corangebg] (m1) at (2.3,1.6) {$\Stk(\mathtt{B}) = \mathtt{c}$};
  \node[hy,draw=corange,fill=corangebg] (m2) at (6.7,1.6) {$\Stk(\mathtt{C}) = \mathtt{b}$};
  \draw[<->,corange,line width=0.9pt] (m1) -- node[above,font=\scriptsize,text=corange] {같은 말} (m2);
  \node[note,anchor=west] at (8.0,1.6) {글자와 짝을 맞바꾼 쌍, 식~\eqref{eq:ac-sym}};
  \node[note,anchor=east] at (-1.1,1.6) {(a)};
  % (b) 작은 봄베의 예
  \node[hy] (h0) at (0,0) {$\Stk(\mathtt{A}) = \mathtt{f}$};
  \node[hy] (h1) at (2.9,0) {$\Stk(\mathtt{B}) = \mathtt{c}$};
  \node[hy,draw=cred,line width=0.9pt] (h2) at (5.8,0) {$\Stk(\mathtt{C}) = \mathtt{d}$};
  \node[hy,draw=cgreen,line width=0.9pt] (h3) at (8.7,0) {$\Stk(\mathtt{A}) = \mathtt{f}$};
  \draw[->] (h0) -- node[lb] {$\Scr_1$} (h1);
  \draw[->] (h1) -- node[lb] {$\Scr_2$} (h2);
  \draw[->] (h2) -- node[lb] {$\Scr_3$} (h3);
  \node[note,text=cgreen,anchor=north] at (8.7,-0.45) {루프는 닫힌다};
  \node[hy,draw=cred,line width=0.9pt,fill=credbg] (h4) at (2.9,-1.7) {$\Stk(\mathtt{C}) = \mathtt{b}$};
  \draw[->,corange,line width=0.9pt] (h1) -- node[right,font=\scriptsize,text=corange] {거울쌍} (h4);
  \draw[cred,dashed,line width=0.8pt] (h4.east) -| (h2.south);
  \node[note,text=cred,anchor=west] at (6.0,-1.3) {\texttt{C}의 짝이 둘 (\texttt{d}, \texttt{b}): 모순};
  \node[note,anchor=east] at (-1.1,0) {(b)};
\end{tikzpicture}
\caption{(a) 플러그보드의 대칭. (b) 작은 봄베의 드럼 위치 1에서 가설 ${\Stk(\mathtt{A}) = \mathtt{f}}$. 루프만 보면 한 바퀴
돌아 \texttt{f}로 돌아오므로(고정점) 살아남는다. 그러나 도중에 나온 $\Stk(\mathtt{B}) = \mathtt{c}$의 거울쌍
$\Stk(\mathtt{C}) = \mathtt{b}$가 이미 나온 $\Stk(\mathtt{C}) = \mathtt{d}$와 부딪친다.}
\label{fig:ac-mirror}
\end{figure}

\Why 그림~\ref{fig:ac-mirror}(b)가 그 효과다. 작은 봄베의 드럼 위치 1에서 루프의 조건을 통과한 것은 \texttt{d}와
\texttt{f} 둘이었다. 그런데 $\Stk(\mathtt{A}) = \mathtt{f}$에서 $\Stk(\mathtt{B}) = \Scr_1(\mathtt{f}) = \mathtt{c}$,
$\Stk(\mathtt{C}) = \Scr_2(\mathtt{c}) = \mathtt{d}$가 나오고, $\Stk(\mathtt{B}) = \mathtt{c}$의 거울쌍은
$\Stk(\mathtt{C}) = \mathtt{b}$다. \texttt{C}의 짝이 둘이니 이 가설은 거짓이다. 루프가 놓친 모순을 거울쌍이 잡았다.
남는 것은 \texttt{d} 하나이고, 이것이 작은 봄베를 만들 때 쓴 실제 플러그보드 \texttt{A}$\leftrightarrow$\texttt{D},
\texttt{B}$\leftrightarrow$\texttt{E}, \texttt{C}$\leftrightarrow$\texttt{F}의 값이다. 참인 가설은 거울쌍을 더해도
지워지지 않는다. $\Stk(\mathtt{A}) = \mathtt{d}$의 귀결 $\Stk(\mathtt{B}) = \mathtt{e}$, $\Stk(\mathtt{C}) = \mathtt{f}$의
거울쌍은 \texttt{D}, \texttt{E}, \texttt{F}에 대한 가설이고 메뉴의 어느 것과도 부딪치지 않는다.

이 사실을 기계에 넣은 것이 웰치먼의 \term{대각선 보드}{diagonal board}다~\cite{welchman1986bomba,carter2010bombe}.
튜링의 처음 설계는 루프가 세 개 이상인 메뉴가 있어야 제대로 쓸 수 있었는데, 대각선 보드를 넣자 루프가 하나뿐이거나
특수한 경우에는 없어도 되었다~\cite{carter2010bombe}. 대각선 보드를 단 첫 기계는 1940년 8월에
가동되었다~\cite{jackson-bh01,cryptomuseum-bombe}. 그림~\ref{fig:g-menu}의 메뉴(루프 1개)를 실제 로터 순서에 걸면,
대각선 보드 없이는 17,576 위치 전부에서 정지하고 대각선 보드가 있으면 2곳에서만 정지한다(보고서 3장).

\Ours \ref{ch:find}장의 그림은 단순하게 하려고 대각선 보드를 빼고 그렸다. 그래서 그 장의 위치 1에서는 \texttt{d}와
\texttt{f}가 함께 꺼진 채 남는다. 대각선 보드를 넣으면 주황 덩어리(\texttt{A}$\to$\texttt{f})가 파란 덩어리와
이어져 \texttt{f}도 켜진다. 대각선 보드를 전선으로 만드는 법은 \ref{sec:c-diagonal}절, Rust 모델에서는
\code{propagate}의 여섯 줄이다(\ref{sec:propagate}절).

\section{정지와 거짓 정지}\label{sec:ac-stop}

이제 봄베가 위치마다 내리는 판정을 적을 수 있다. 테스트 글자 $T$를 하나 고르고, $T$에 대한 가설 하나를 넣어 귀결을
모두 모은다. 귀결이 $T$의 가설 26개를 \emph{모두} 덮으면 이 위치를 버린다. 하나라도 남으면 기계가 멈춘다. 이것이
\term{정지}{stop}다. 1차 자료는 정지 조건을 ``입력에 통전되지 않은 점이 있는 것''으로 적는다~\cite[32]{us6812report1944}.
테스트 글자로는 선이 가장 많이 모이는 글자를 골랐다~\cite[7]{us6812report1944}.

\Why 이 판정은 한쪽으로만 믿을 수 있다.

\begin{itemize}
\item \textbf{맞는 위치는 버려지지 않는다.} 맞는 로터 순서와 코어 위치에서는 모든 추론이 참이다. 참인 가설에서는 참인
것만 나온다. 그러면 거짓인 가설에서 참인 가설이 나올 수도 없다. 그런 길이 있다면 추론을 거꾸로 따라 참에서 거짓이
나오기 때문이다. 그러므로 $T$의 참인 짝은 어떤 거짓 가설의 귀결에도 들지 않고, 26개가 다 덮이는 일은 없다.
\item \textbf{멈췄다고 맞는 위치는 아니다.} 틀린 위치에서도 모순이 우연히 드러나지 않으면 멈춘다. 이것이
\term{거짓 정지}{false stop}다. 정지는 답이 아니라 \emph{후보}다.
\end{itemize}

그림~\ref{fig:ac-register}는 작은 봄베에서 테스트 글자 \texttt{A}의 가설 여섯 개가 어떻게 되었는지를 칸으로 그린 것이다.
칠한 칸이 지워진 가설이다. 이 칸 줄을 기계에서는 \term{테스트 레지스터}{test register}라 부른다(\ref{sec:f-register}절).

\begin{figure}[ht]
\centering
\begin{tikzpicture}[x=6mm,y=-8mm,
    on/.style={draw=white,fill=cblue,minimum size=5.6mm,inner sep=0pt,font=\ttfamily\small,text=white},
    off/.style={draw=cred,line width=0.9pt,fill=white,minimum size=5.6mm,inner sep=0pt,font=\ttfamily\small,text=cred}]
  \foreach \r in {0,1,2} {
    \foreach \l [count=\i from 0] in {a,b,c,d,e,f} { \node[on] at (\i,\r) {\l}; }
  }
  \node[off] at (3,0) {d}; \node[off] at (5,0) {f};
  \node[off] at (3,1) {d};
  \node[anchor=east,font=\footnotesize] at (-0.8,0) {드럼 위치 1, 대각선 보드 없음};
  \node[anchor=east,font=\footnotesize] at (-0.8,1) {드럼 위치 1, 대각선 보드 있음};
  \node[anchor=east,font=\footnotesize] at (-0.8,2) {드럼 위치 2};
  \node[anchor=west,font=\footnotesize] at (5.8,0) {정지: 후보 \texttt{d}, \texttt{f}};
  \node[anchor=west,font=\footnotesize] at (5.8,1) {정지: 후보 \texttt{d}};
  \node[anchor=west,font=\footnotesize] at (5.8,2) {다음 위치로};
\end{tikzpicture}
\caption{작은 봄베에서 가설 $\Stk(\mathtt{A}) = \mathtt{a}$를 넣었을 때 테스트 글자 \texttt{A}의 여섯 가설. 파란 칸은
지워진 가설, 빨간 테두리 칸은 남은 가설이다. 맨 아래 줄처럼 여섯 칸이 다 칠해지면 그 위치를 버린다.}
\label{fig:ac-register}
\end{figure}

\paragraph{거짓 정지는 얼마나 나오는가.} 대각선 보드가 없을 때는 앞 절들의 논리로 어림할 수 있다. 메뉴의 덩어리에
루프가 $L$개 있다고 하자.

\begin{itemize}
\item $L = 0$이면 추론이 테스트 글자로 돌아오지 않는다. 가설 하나에서 $T$의 가설은 그 자신 하나만 나오므로 모든 위치에서
정지한다.
\item $L = 1$이면 귀결은 넣은 가설이 속한 $P$의 순환이다. 26칸이 다 덮이려면 $P$가 26글자를 한 순환으로 돌아야 한다.
무작위 순열이 그럴 확률은 $25!/26! = 1/26$이다. 그러니 위치의 약 $25/26$, 곧 약 16,900곳에서 정지한다.
\item $L \geq 2$이면 거짓 정지는 대개 모든 루프에 공통인 고정점이 있을 때 생긴다. 글자 하나가 $L$개 루프를 모두 통과할
확률이 약 $(1/26)^L$이고 글자가 26개이므로, 로터 순서 하나의 17,576 위치에서
\begin{equation}
  \text{거짓 정지} \approx 17{,}576 \times 26 \times \left(\frac{1}{26}\right)^{L} = 17{,}576 \times 26^{\,1-L}. \label{eq:ac-falsestop}
\end{equation}
\end{itemize}

표~\ref{tab:ac-estimate}가 이 어림과 보고서 7장의 측정을 나란히 놓은 것이다. 루프 하나가 정지를 약 26분의 1로
줄인다는 것이 그대로 보인다. 대각선 보드가 없으면 루프가 세 개는 있어야 정지가 수십 개로 줄어든다. 튜링의 처음
설계에 루프가 세 개 이상 필요했다는 기록~\cite{carter2010bombe}과 맞는 산수다.

\begin{table}[ht]
\centering
\caption{로터 순서 하나(17,576 위치)에서의 거짓 정지 수. 측정은 글자 12개의 무작위 메뉴를 경우마다 400번 돌린
평균이다(보고서 7장). 어림은 이 절의 계산이다.}
\label{tab:ac-estimate}
\small
\begin{tabular}{c r r r}
\toprule
루프 수 $L$ & 어림(대각선 보드 없음) & 측정(대각선 보드 없음) & 측정(대각선 보드 있음) \\
\midrule
0 & 17,576 & 17,576 & 673 \\
1 & 약 16,900 & 16,960 & 25.3 \\
2 & 676 & 722 & 0.88 \\
3 & 26 & 26.6 & 0.03 \\
\bottomrule
\end{tabular}
\end{table}

대각선 보드가 있으면 어림이 복잡해진다. 거울쌍이 메뉴의 글자 수만큼 조건을 더하기 때문이다. 표의 오른쪽 열에서 보듯
정지는 크게 줄고, 루프 하나마다 약 26분의 1이 되는 모양은 그대로다(보고서 7장). 그림~\ref{fig:g-menu}의 메뉴
(글자 14개, 루프 1개)에 대한 튜링의 추정은 로터 순서당 약 1.6회이고~\cite{carter-stops,wikipedia-bombe}, 이
프로젝트에서 Enigma~I의 120가지 로터 순서를 돌리면 정지가 157회, 로터 순서당 1.31회 나온다. 23글자 크립의 메뉴(루프
4개)로는 2회로 준다(보고서 7장). 실제 운용에서는 로터 순서당 정지가 한 번 이상(기계가 제대로 돌았는지 확인하려고),
네 번쯤 이하(한 사람이 체킹을 감당할 수 있는 만큼)가 되도록 메뉴를 짰다~\cite[31, 44, 49]{us6812report1944}.

\Ours 정지 판정은 Rust 모델에서 \code{if reg != FULL} 한 줄이다(\ref{ch:rust-bombe}장). 세 기계에서 이 판정이
무엇으로 만들어지는지는 \ref{sec:f-register}절에서 본다.

\section{체킹: 정지에서 키까지}\label{sec:ac-check}

정지 하나가 알려 주는 것은 로터 순서, 코어 위치, 테스트 글자의 짝 후보뿐이다. 이것을 검증하고 나머지 키를 구하는
뒷단계가 \term{체킹}{checking}이다.

\Why 체킹은 봄베와 같은 추론을 반대 방향으로 쓴다. 봄베는 거짓 가설 하나로 나머지를 지웠다. 체킹은 남은 후보를
참이라고 놓고 식~\eqref{eq:ac-moved}로 메뉴의 선을 하나씩 따라가 귀결을 \emph{읽는다}. 후보가 참이라면 글자마다 짝이
하나씩만 나와야 하고, 루프는 한 바퀴 돌아 처음 값으로 닫혀야 하며, 거울쌍이 서로 맞아야 한다. 작은 봄베에서는
$\Stk(\mathtt{A}) = \mathtt{d}$에서 $\Stk(\mathtt{B}) = \mathtt{e}$, $\Stk(\mathtt{C}) = \mathtt{f}$가 나오고, 루프가
\texttt{d}로 닫힌다(그림~\ref{fig:ac-check}). 플러그보드 세 쌍 \texttt{A}$\leftrightarrow$\texttt{D},
\texttt{B}$\leftrightarrow$\texttt{E}, \texttt{C}$\leftrightarrow$\texttt{F}가 읽혔다.

\begin{figure}[ht]
\centering
\begin{tikzpicture}[x=1cm,y=-5.5mm,
    dot/.style={circle,draw=cgray!70,fill=white,inner sep=1.3pt}]
  \foreach \col/\name/\key in {0/A/A,3/B/B,6/C/C,9/A/Z} {
    \node[font=\small\bfseries] at (\col,-1.1) {\texttt{\name}};
    \foreach \l [count=\i from 0] in {a,b,c,d,e,f} { \node[dot] (\key\l) at (\col,\i) {}; }
  }
  \foreach \l [count=\i from 0] in {a,b,c,d,e,f} {
    \node[font=\ttfamily\scriptsize,anchor=east] at (-0.15,\i) {\l};
    \node[font=\ttfamily\scriptsize,anchor=west] at (9.15,\i) {\l};
  }
  \draw[line width=1.3pt,cgreen] (Ad) -- node[above=1.5pt,sloped,font=\scriptsize] {$\Scr_1$} (Be)
     -- node[above=1.5pt,sloped,font=\scriptsize] {$\Scr_2$} (Cf) -- node[above=1.5pt,sloped,font=\scriptsize] {$\Scr_3$} (Zd);
  \foreach \n in {Ad,Be,Cf,Zd} { \node[circle,fill=cgreen,inner sep=1.8pt] at (\n) {}; }
  \node[note,anchor=north] at (1.5,6.2) {$\Stk(\mathtt{A})=\mathtt{d}$이면 $\Stk(\mathtt{B})=\mathtt{e}$};
  \node[note,anchor=north] at (7.5,6.2) {루프가 \texttt{d}로 닫힌다};
\end{tikzpicture}
\caption{작은 봄베의 정지(드럼 위치 1)를 체킹하는 것. 남은 후보 $\Stk(\mathtt{A}) = \mathtt{d}$ 한 가닥만 루프를 따라
간다. 오른쪽 끝 열은 다시 \texttt{A}다. 글자마다 값이 하나씩 나오고 처음 값으로 돌아오므로 통과다.}
\label{fig:ac-check}
\end{figure}

실제 체킹에서는 귀결이 메뉴 밖의 글자로 나가기도 한다. 그래서 판정이 세 가지로 나뉘었다(표~\ref{tab:ac-verdicts}).

\begin{table}[ht]
\centering
\caption{체킹 결과의 분류~\cite[12--16, 49--50]{us6812report1944}(보고서 3장의 표를 줄임).}
\label{tab:ac-verdicts}
\small
\begin{tabularx}{\linewidth}{l>{\raggedright\arraybackslash}X>{\raggedright\arraybackslash}X}
\toprule
용어 & 뜻 & 처리 \\
\midrule
confirmation & 메뉴 글자 $X$의 짝이 메뉴 글자 $Y$이고 $Y$의 짝이 $X$ & 좋은 징후 \\ \addlinespace
legal contradiction & 메뉴 글자 둘이 메뉴 밖의 같은 글자를 짝으로 가짐 & 기계로서는 정상인 정지(대각선 보드가 거르지 못하는 모순) \\ \addlinespace
illegal contradiction & 메뉴 글자끼리의 짝이 비대칭이거나 루프가 닫히지 않음 & 기계 고장. 대각선 보드가 있으면 나올 수 없다 \\
\bottomrule
\end{tabularx}
\end{table}

메뉴를 통과한 정지는 남은 정보로 더 시험한다.

\begin{enumerate}
\item \textbf{크립 전체.} 메뉴에 넣지 못한 크립의 글자(그림~\ref{fig:g-menu}에서 뺀 \texttt{O}--\texttt{F},
\texttt{S}--\texttt{N}, 선을 고르며 버린 것)도 모두 식~\eqref{eq:ac-moved}를 만족해야 한다. 짝을 아는 글자에서 모르는
글자로 짝을 채워 가다 모순이 나오면 버린다. 예제의 정지 157개 가운데 이 단계까지 통과하는 것은 하나다(\ref{ch:rust-check}장).
\item \textbf{링 설정.} 봄베가 찾은 것은 코어 위치다. 창 글자와 링 설정은 따로 정해지지 않는다. 둘을 가르는 것은 로터가
옆 로터를 언제 넘기는가뿐이므로, 넘어감이 들어 있는 긴 전문을 복호화해 보아야 정해진다.
\item \textbf{남은 스테커.} 크립에 나오지 않은 글자의 짝은 논리로 정할 수 없다. 전문 전체를 복호화해 언어다운 정도를
보며 고른다. 이것은 추론이 아니라 추측이다.
\end{enumerate}

그 끝에 평문이 나온다. 봄베를 운용한 출장소의 체킹은 스테커의 일관성 검사까지였고, 메뉴에 없는 글자의 스테커와 실제 링
설정은 Hut~6에서 손으로 구했다~\cite{ellsbury1998bombe,carter2010bombe}. 미 해군의 4로터 봄베는 정지한 위치에서 스테커를
스스로 검사해 통과한 것만 인쇄했다~\cite{ekhall-hallenberg2018}.

\Ours 위의 세 단계를 포함한 체킹 전체가 \code{check\_stop}이다(\ref{ch:rust-check}장). 그 장의 1단계가
그림~\ref{fig:ac-check}, 2단계가 위의 1번, 3--4단계가 위의 2--3번이다.

\section{넘어감과 4로터}\label{sec:ac-turnover}

\paragraph{넘어감이 없다는 가정.} 지금까지 $\Scr_k$가 ``위치만 정하면 계산된다''고 했다. 이것은 크립 구간에서 오른쪽
로터만 움직인다는 가정 위에 서 있다. 그러면 모든 $\Scr_k$가 가운데와 왼쪽 로터의 위치를 함께 쓰고, 오른쪽 로터만
$k$칸씩 어긋난다. 위치 하나가 $\Scr_1, \Scr_2, \ldots$를 모두 정하므로 로터 순서마다 $26^3$ 위치만 훑으면 된다. 봄베의
드럼에는 노치가 없고 스크램블러 사이의 상대 위치는 고정이다. 곧 봄베는 메뉴 구간 안에서 가운데 로터가 움직이지 않는다고
가정한다~\cite[4]{us6812report1944}.

크립 위치 $t$에서 가운데 로터가 넘어가면 $k \geq t$인 선의 스크램블러는 가운데 로터가 한 칸 더 간 것이다
(그림~\ref{fig:ac-turnover}). 이 선들에 원래의 $\Scr_k$를 걸면 맞는 위치에서도 식~\eqref{eq:ac-moved}가 틀리고, 정지가
나오지 않는다. 오른쪽 로터는 26글자에 한 번 가운데 로터를 넘기므로 크립이 길수록 위험이 크다. 문헌은 23글자 크립의
구간에 넘어감이 들어 있을 확률을 $23/26$으로 적는다~\cite{ellsbury1998bombe}. 크립을 나누거나, 뒤쪽 선의 가운데
드럼을 한 칸 넘긴 메뉴를 따로 돌려 대처했다~\cite{ellsbury1998bombe}\cite[37]{us6812report1944}.

\begin{figure}[ht]
\centering
\begin{tikzpicture}[x=1.05cm,y=1cm,
    cell/.style={draw=cgray,minimum width=9.5mm,minimum height=6mm,inner sep=0pt,font=\scriptsize}]
  \node[anchor=east,font=\footnotesize] at (0.4,0.75) {크립 위치 $k$};
  \node[anchor=east,font=\footnotesize] at (0.4,0) {오른쪽 로터};
  \node[anchor=east,font=\footnotesize] at (0.4,-0.75) {가운데 로터};
  \foreach \n in {1,...,10} {
    \node[font=\scriptsize,text=cgray] at (\n,0.75) {\n};
    \node[cell,fill=cbluebg] at (\n,0) {$r{+}\n$};
    \ifnum\n<6
      \node[cell,fill=white] at (\n,-0.75) {$m$};
    \else
      \node[cell,fill=corangebg,draw=corange] at (\n,-0.75) {$m{+}1$};
    \fi
  }
  \draw[corange,dashed,line width=0.9pt] (5.5,1.05) -- (5.5,-1.15) node[below,font=\scriptsize,text=corange] {넘어감 $t = 6$};
\end{tikzpicture}
\caption{넘어감 가설. 위치 하나 $(m, r)$이 모든 $\Scr_k$를 정하되, $k \geq t$인 선은 가운데 로터가 한 칸 더 간
스크램블러를 쓴다. 가설이 없는 봄베는 아래 줄이 모두 $m$이라고 가정한다.}
\label{fig:ac-turnover}
\end{figure}

\Ours Rust 모델은 ``크립 위치 $t$에서 넘어갔다''는 가설을 받아 $k \geq t$인 선의 표를 두 번째 벌에서 고른다
(\ref{ch:rust-bombe}장의 넘어감 가설 절). 가설마다 구동이 한 번씩 늘어난다(\ref{sec:s-count}절).

\paragraph{4로터.} 해군의 M4에는 네 번째 로터가 더 있다. 이 로터는 전문을 치는 동안 움직이지 않으므로, 얇은 반사판과
묶으면 ``반사판 26종 가운데 하나''로 볼 수 있다~\cite{cryptomuseum-bombe}. $\Scr_k$가 고정점 없는 되돌릴 수 있는 치환이라는
것은 그대로이고, 이 장의 논리도 그대로다. 위치만 로터 순서당 $26^4 = 456{,}976$가지로 26배가 된다. 로터 순서 1,344가지를
모두 하면 614,175,744 위치다(\ref{sec:s-count}절). 이 프로젝트의 Rust 모델로 M4 전체를 돌리면 정지가 176회 나오고 그중
체킹을 통과한 해가 하나다(보고서 7장).

\section{폴란드의 방법}\label{sec:ac-polish}

튜링의 봄베보다 먼저, 폴란드 암호국은 다른 길로 에니그마를 풀었다. 크립 대신 \emph{전문 키를 보내는 절차}를 썼다.
이 절은 그 수학을 짧게 본다. 출처는 레예프스키가 직접 쓴 논문이다~\cite{rejewski1980application}.

\paragraph{두 번 암호화한 전문 키.} 1938년 9월 15일까지 독일군의 절차는 이랬다~\cite{rejewski1980application}. 운용자는
그날 정해진 기본 위치에 로터를 맞추고, 전문마다 세 글자의 전문 키를 골라 \emph{두 번} 쳤다. 그렇게 나온 여섯 글자가
전문 맨 앞에 붙었다. 그날의 모든 전문이 같은 기본 위치에서 키를 암호화했고, 첫째 글자와 넷째 글자는 같은 글자를
암호화한 것이다.

기본 위치에서 $i$번째 글자를 치는 에니그마를 $E_i = \Stk\,\Scr_i\,\Stk$로 적자. 여기서 $\Scr_i$는 크립 위치가 아니라
기본 위치에서 $i$번째 글자를 칠 때의 스크램블러다. 전문 키의 첫 글자를 $x$라 하면 암호문의
첫째 글자는 $a = E_1(x)$, 넷째 글자는 $b = E_4(x)$다. $x$는 모르지만 $E_1$은 되돌릴 수 있으므로 $x = E_1(a)$이고,
\begin{equation}
  b = E_4 E_1 (a) \label{eq:ac-indicator}
\end{equation}
다. 첫째 글자에서 넷째 글자로 가는 대응은 치환 $E_4E_1$이다. 그날 전문이 충분히 모이면(약 80개) 이 치환 전체가
드러난다. 둘째--다섯째, 셋째--여섯째도 같다~\cite{rejewski1980application}(그림~\ref{fig:ac-indicator}).

\begin{figure}[ht]
\centering
\begin{tikzpicture}[x=1cm,y=1cm,
    bx/.style={draw=cgray,minimum width=6.5mm,minimum height=6.5mm,inner sep=0pt,font=\ttfamily\small}]
  % 여섯 글자
  \foreach \c [count=\i] in {d,m,q,v,b,n} { \node[bx] (t\i) at (\i*0.75,0) {\c}; }
  \draw[cblue,line width=0.8pt] (t1.north) to[bend left=50] node[above,font=\scriptsize] {$E_4E_1$} (t4.north);
  \draw[cgray!80] (t2.south) to[bend right=35] (t5.south);
  \draw[cgray!80,dashed] (t3.south) to[bend right=60] (t6.south);
  \node[note,anchor=north] at (2.625,-1.2) {전문 하나의 머리 여섯 글자};
  % 대응
  \begin{scope}[xshift=7.2cm]
    \node[ltr] (a) at (0,0) {d};
    \node[ltr,fill=corangebg,draw=corange] (x) at (2,0) {$x$};
    \node[ltr] (b) at (4,0) {v};
    \draw[<->] (a) -- node[above,font=\scriptsize] {$E_1$} (x);
    \draw[->] (x) -- node[above,font=\scriptsize] {$E_4$} (b);
    \draw[->,cblue,line width=1pt] (a) to[bend right=35] node[below,font=\scriptsize] {$E_4E_1$} (b);
    \node[note,anchor=south] at (2,0.45) {키의 첫 글자(모름)};
  \end{scope}
\end{tikzpicture}
\caption{두 번 암호화한 전문 키. 레예프스키의 논문에 나오는 예 \texttt{dmq vbn}에서 첫째 글자 \texttt{d}와 넷째 글자
\texttt{v}는 같은 글자 $x$를 암호화한 것이다. 그러므로 $E_4E_1$은 \texttt{d}를 \texttt{v}로 보낸다. 논문은 이 치환을
왼쪽부터 곱하는 표기로 $AD$라 적는다.}
\label{fig:ac-indicator}
\end{figure}

\paragraph{레예프스키의 정리.} $E_1$과 $E_4$는 둘 다 쌍을 맞바꾸기만 하는 치환이다(고정점 없는 대합,
\ref{sec:am-properties}절). 레예프스키는 이런 치환 둘의 곱에 관한 정리를 썼다~\cite{rejewski1980application}.

\begin{keybox}[title={레예프스키의 정리}]
$X$와 $Y$가 둘 다 서로소인 맞바꿈(쌍)만으로 이루어진 치환이면, 곱 $YX$의 순환은 길이마다 짝수 개씩 있다. 같은 길이의
순환이 짝을 지어 나온다.
\end{keybox}

까닭은 그림으로 보인다(그림~\ref{fig:ac-rejewski}). 글자를 점으로 놓고 $X$의 쌍을 실선으로, $Y$의 쌍을 점선으로 잇는다.
모든 점에 실선 하나와 점선 하나가 닿으므로 이 그림은 실선과 점선이 번갈아 나오는 닫힌 다각형들로 나뉜다. 꼭짓점이
$2m$개인 다각형 하나에서 $YX$는 실선 하나, 점선 하나를 건너 두 칸씩 움직인다. 그러면 홀수 번째 꼭짓점끼리 길이 $m$인
순환 하나, 짝수 번째 꼭짓점끼리 길이 $m$인 순환 하나가 생긴다. 다각형마다 같은 길이의 순환이 둘씩 나온다. ($X$와 $Y$에
같은 쌍이 있으면 꼭짓점 둘짜리 다각형이 되어 길이 1인 순환, 곧 고정점 둘이 나온다.)

\begin{figure}[ht]
\centering
\begin{tikzpicture}[x=1cm,y=1cm,
    od/.style={ltr,fill=cbluebg,draw=cblue},
    ev/.style={ltr,fill=cgreenbg,draw=cgreen},
    xe/.style={line width=1.1pt,corange},
    ye/.style={line width=1.1pt,cpurple,dashed}]
  % 다각형
  \node[od] (a) at (90:1.45) {a};
  \node[ev] (b) at (30:1.45) {b};
  \node[od] (e) at (-30:1.45) {e};
  \node[ev] (d) at (-90:1.45) {d};
  \node[od] (f) at (-150:1.45) {f};
  \node[ev] (c) at (150:1.45) {c};
  \draw[xe] (a) -- (b); \draw[xe] (e) -- (d); \draw[xe] (f) -- (c);
  \draw[ye] (b) -- (e); \draw[ye] (d) -- (f); \draw[ye] (c) -- (a);
  \node[note,anchor=north] at (0,-2.0) {실선 $X$, 점선 $Y$};
  % 곱의 순환
  \begin{scope}[xshift=4.6cm]
    \node[od] (pa) at (0,1.0) {a};
    \node[od] (pe) at (0.85,-0.45) {e};
    \node[od] (pf) at (-0.85,-0.45) {f};
    \draw[->,cblue] (pa) -- (pe); \draw[->,cblue] (pe) -- (pf); \draw[->,cblue] (pf) -- (pa);
  \end{scope}
  \begin{scope}[xshift=7.6cm]
    \node[ev] (pb) at (0,1.0) {b};
    \node[ev] (pc) at (0.85,-0.45) {c};
    \node[ev] (pd) at (-0.85,-0.45) {d};
    \draw[->,cgreen] (pb) -- (pc); \draw[->,cgreen] (pc) -- (pd); \draw[->,cgreen] (pd) -- (pb);
  \end{scope}
  \node[note,anchor=north] at (6.1,-2.0) {$YX = (\mathtt{a}\,\mathtt{e}\,\mathtt{f})(\mathtt{b}\,\mathtt{c}\,\mathtt{d})$};
\end{tikzpicture}
\caption{레예프스키의 정리를 작은 봄베의 스크램블러로 본 것. $X = \Scr_1 = (\mathtt{ab})(\mathtt{cf})(\mathtt{de})$,
$Y = \Scr_3 = (\mathtt{ac})(\mathtt{be})(\mathtt{df})$. 왼쪽 육각형을 따라 두 칸씩 가면 파란 꼭짓점끼리, 초록
꼭짓점끼리 길이 3인 순환이 하나씩 생긴다(오른쪽). 예를 들어 \texttt{a}는 $X$로 \texttt{b}, $Y$로 \texttt{e}에 간다.}
\label{fig:ac-rejewski}
\end{figure}

\paragraph{플러그보드가 사라진다.} 이 곱에서 플러그보드가 어떻게 되는지 보자. $\Stk\Stk$는 아무것도 바꾸지 않으므로
\begin{equation}
  E_4E_1 = \Stk\,\Scr_4\,\Stk\;\Stk\,\Scr_1\,\Stk = \Stk\,(\Scr_4\Scr_1)\,\Stk \label{eq:ac-conj}
\end{equation}
다. 치환 $Q$가 $u$를 $v$로 보내면 $\Stk Q \Stk$는 $\Stk(u)$를 $\Stk(v)$로 보낸다. 순환의 글자 이름만 $\Stk$로 바뀌고
순환의 길이는 그대로다. 그러므로 $E_4E_1$의 순환 길이들은 플러그보드와 무관하고, 로터 순서와 기본 위치만으로 정해진다.
레예프스키는 그날의 세 곱을 묶어 그날의 \term{특성}{characteristic}이라 불렀다. 폴란드 팀은 \term{사이클로미터}{cyclometer}라는
장치로 로터 순서 6가지 $\times$ 위치 17,576가지 $= 105{,}456$가지에 대해 특성의 순환 길이와 개수를 목록으로 만들었다.
그날의 특성과 같은 모양을 목록에서 찾으면 로터 순서와 플러그보드를 얻었고, 나머지 키는 다른 방법으로
구했다~\cite{rejewski1980application}.

봄베와 같은 착상이 다른 모양으로 들어 있다. 식~\eqref{eq:ac-moved}는 $\Stk$를 한쪽으로 넘겨 ``모르는 것''으로 따로 두었고,
식~\eqref{eq:ac-conj}는 $\Stk$가 순환의 모양에서 지워지게 했다. 어느 쪽이든 큰 쪽 키(플러그보드)를 대입하지 않고 작은 쪽
키(로터)만 훑는다.

\paragraph{절차가 바뀌자.} 1938년 9월 15일부터 SD(친위대 보안국)를 뺀 모든 부대에서 운용자는 전문마다 기본 위치를
스스로 골라 전문 머리에 평문으로 적게 되었다. 전문 키는 여전히 두 번 암호화했다. 그날 공통의 곱 $E_4E_1$이 없어졌으므로 특성 목록은 쓸모를 잃었다. 남은
실마리는 첫째와 넷째 글자가 \emph{같은} 전문, 곧 곱의 고정점이었고, 고정점이 있고 없음도 플러그보드와 무관하다.
지갈스키의 구멍 뚫은 종이가 이것을 썼다. 1938년 12월 15일에는 로터가 세 개에서 다섯 개 가운데 셋을 고르는 것으로
늘었고, 플러그보드의 쌍도 6쌍에서 13쌍까지 점점 늘었다~\cite{rejewski1980application}. 폴란드 팀은 1938년에 전기기계식 장치
봄바(bomba)도 만들었는데, 이것 역시 당시의 지시자 절차를 이용하는 장치였고 튜링--웰치먼의 봄베와는 전혀 다른
기계다~\cite{cryptomuseum-bombe,jackson-bh03}. 1939년 7월 폴란드는 이 성과를 영국과 프랑스에
넘겼다~\cite{rejewski1980application,gchq-pyry2020}.

지시자 절차에 기대는 방법은 절차가 바뀌면 쓸 수 없다. 튜링은 절차가 아니라 전문의 본문, 곧 크립에 기대는 기계를
설계했다~\cite{jackson-bh06,ellsbury1998bombe,hodges-biography}. 이 장의 앞부분이 그 기계의 논리다.

\section{정리}

\begin{itemize}
\item 키의 대부분은 플러그보드다. 봄베는 로터 순서와 코어 위치만 훑고, 플러그보드는 논리로 거른다.
\item 크립은 ``어떤 글자도 자기 자신이 되지 않는다''로 놓을 자리를 추린다. 틀린 자리가 살아남을 확률은 글자마다 $25/26$씩 준다.
\item 플러그보드를 한쪽으로 넘기면 $\Stk(c) = \Scr_k(\Stk(p))$(식~\eqref{eq:ac-moved})다. 가설 하나가 다른 가설 하나를
계산해 내고, 이 추론은 양쪽으로 읽힌다. 메뉴는 이 추론들의 그래프다.
\item 루프를 따라가면 $\Stk(X_0)$는 스크램블러 곱 $P$의 고정점이어야 한다(식~\eqref{eq:ac-loop}). 가설 하나를 넣고 귀결을
따라가면 $P$의 순환이 나오고, 순환의 가설은 함께 넘어진다. 루프 하나가 후보를 약 26분의 1로 줄인다.
\item 거울쌍 $\Stk(X)=y \iff \Stk(Y)=x$(대각선 보드)는 루프 없이도 추론을 되돌려 보내, 필요한 루프의 수를 줄인다.
\item 테스트 글자의 가설 26개가 모두 지워지면 위치를 버리고, 아니면 정지한다. 맞는 위치는 버려지지 않지만, 정지는 후보일
뿐이다. 체킹이 남은 후보를 따라가 키를 읽고 평문을 낸다.
\item 이 논리는 크립 구간에 넘어감이 없다는 가정 위에 있고, 4로터에서도 그대로다. 폴란드의 방법은 같은 약점을 지시자
절차로 이용했다.
\end{itemize}

이 장의 추론은 모두 ``가설 하나에서 가설 하나로''의 사슬이었다. 가설은 676개이고 추론은 양쪽으로 읽힌다. 가설을 전선
한 가닥으로, 추론을 전선의 연결로 바꾸면, 귀결을 따라가는 일을 전류가 한순간에 해 준다. 이 추론을 전선으로 만든 것이
봄베다. 다음 장이 그 이야기다.
